\documentclass[10pt,a4paper]{amsart}
\usepackage[margin=19mm]{geometry}
\usepackage{amsmath,amssymb,mathtools}
\usepackage[scr=rsfs,cal=euler]{mathalfa}
\usepackage{xfrac}
\usepackage[unicode,hidelinks,bookmarks=false]{hyperref}
\newcommand{\ie}{i.e.\ }
\newcommand{\cf}{cf.\ }
\newcommand{\resp}{respectively\ }
\newcommand{\Subsection}{\subsection}
\providecommand{\nobreakdash}{\nobreak\hbox{-}}
\setlength{\emergencystretch}{2em}
\multlinegap0pt
\usepackage{bm}
\usepackage{bbm}
\usepackage{dsfont}
\usepackage{xspace}
\usepackage{cancel}
\usepackage{mathtools}
\usepackage[shortlabels]{enumitem}
\usepackage{amsthm}
\makeatletter
\@ifundefined{theorem}{\newtheorem{theorem}{Theorem}}{}
\@ifundefined{lemma}{\newtheorem{lemma}[theorem]{Lemma}}{}
\@ifundefined{observation}{\newtheorem{observation}[theorem]{Observation}}{}
\@ifundefined{proposition}{\newtheorem{proposition}[theorem]{Proposition}}{}
\@ifundefined{remark}{\newtheorem{remark}[theorem]{Remark}}{}
\makeatother
\DeclareMathOperator{\link}{\mathrm{link}}
\title[Cone and constrained colorful Carath\'eodory Theorems]{Cone and constrained colorful Carath\'eodory Theorems}
\author{Pavle V. M. Blagojevi\'c, Roman N. Karasev, B\'alint Zsigri}
\date{Source: arXiv:2607.06172v1 (CC BY 4.0)}
\begin{document}
\maketitle

\noindent\emph{Source and licence.} Cone and constrained colorful Carath\'eodory Theorems. Version arXiv:2607.06172v1 (CC BY 4.0), distributed under the Creative Commons Attribution 4.0 International licence, as recorded in the arXiv licence metadata for this exact version and shown on the abstract page https://arxiv.org/abs/2607.06172v1. Full rights notice: licenses/arxiv-2607.06172-CC-BY-4.0.txt.\par\smallskip

\noindent\textit{Editorial scope.} Counted unit: the Theorem whose source label is \texttt{thm:nerve-closed-cellular} in the article (source0.tex lines 1443--1445, source label \texttt{thm:nerve-closed-cellular}), delivered together with the article's own complete original proof of it (source0.tex lines 1446--1469, one \texttt{proof} environment of 24 lines). No mathematics is written, repaired or re-derived: statement and proof are the article's own text line for line. Required context retained uncounted: def:om:elim (lines 983--985); lem:disks-in-oms (lines 1096--1101); obs:pseudospheres-regular-cw (lines 1102--1104); rmk:reorientation (lines 1128--1130); obs:curly-vs-print (lines 1132--1136); prop:htpy-type-zero-avoiding (lines 1222--1243); prop:htpy-type-zero-avoiding:nerve (lines 1224--1227); prop:htpy-type-elem-avoiding (lines 1408--1420); prop:htpy-type-elem-avoiding:nerve (lines 1410--1413); prop:htpy-type-elem-avoiding:void (lines 1410--1413); prop:htpy-type-elem-avoiding:nonpointed (lines 1415--1418); prop:htpy-type-elem-avoiding:pointed (lines 1415--1418). Every definition, display and label of those lines is retained unchanged. Omitted from this package and not redistributed: 1--982, 986--1095, 1105--1127, 1131--1131, 1137--1221, 1228--1407, 1414--1414, 1419--1442, 1470--1735. Nothing inside a retained excerpt is omitted or altered: every retained body is delivered exactly as the article's lines are written, so no omission receipt of any kind accompanies this package. Citations: the retained excerpts carry no citation command at all, so no bibliography entry of the article is transcribed and no reference is redistributed. Reference adaptation: none was needed and none was made; every reference inside the delivered unit resolves inside it, so no unresolved reference or citation is produced. Printed slips kept literal: no printed word, symbol, hypothesis or number of the counted statement or of its proof was altered. Layout and class: the article's own class is \texttt{amsart} with packages including anysize, color, inputenc, xcolor, hyperref, float, amsfonts, amssymb, amsmath, amsthm, url, paralist, tikz, rotating; the wrapper is the lane's pinned \texttt{amsart} editorial wrapper at 10pt on A4, which restates the theorem-like environments the retained text needs and adds the article's own macro definitions that the retained text needs (\texttt{\textbackslash link}), with extra wrapper packages (bm, bbm, dsfont, xspace, cancel, mathtools, enumitem). The delivered numbering is the unit's own. Assets: no retained span contains a figure environment, a graphics-inclusion command, a PDF-page inclusion, a table or a TikZ picture; no image file is copied into this package, the package declares no asset dependency and no image file is redistributed with it. Licence: version arXiv:2607.06172v1 of this article is distributed under the Creative Commons Attribution 4.0 International licence, as recorded in the arXiv licence metadata for this exact version and shown on the abstract page https://arxiv.org/abs/2607.06172v1. A full rights notice is delivered at licenses/arxiv-2607.06172-CC-BY-4.0.txt. Characters: every retained excerpt is pure ASCII, so the delivered engine, XeLaTeX with its default Latin Modern text font, reports no missing character.
    \begin{enumerate}[resume,label=\normalfont{(V\arabic*')}]
        \item\label{def:om:elim} For all $\Pi,\Xi\in\mathscr{L}(\mathscr{O})$ and $v\in\Pi^+\cap\Xi^-$ there is a $\Theta\in\mathscr{L}(\mathscr{O})$ such that $\Theta(v)=0$ and if $w\in V-S(\Pi,\Xi)$ then $\Theta(w)=\Pi(w)\circ\Xi(w)$.
    \end{enumerate}

\begin{lemma}\label{lem:disks-in-oms}
    Let $\mathfrak{A}$ be a signed arrangement of pseudospheres, $\mathscr{O}$ the associated oriented matroid, and $\Phi\in\mathscr{L}(\mathscr{O})-{0}$. Then  $D^-(\Phi^-)\cap S(\Phi^0)\cap D^+(\Phi^+)=\mathfrak{s}_\mathfrak{A}^{-1}(\{\Psi\in\mathscr{L}(\mathscr{O}):\Psi\leq\Phi\})$ is a $\dim S(\Phi^0)$-dimensional disk. Furthermore, its relative boundary as a subset of the sphere $S(\Phi^0)$ is 
    \[
        D^-(\Phi^-)\cap S(\Phi^0)\cap D^+(\Phi^+)-B^-(\Phi^-)\cap S(\Phi^0)\cap B^+(\Phi^+)=\mathfrak{s}_\mathfrak{A}^{-1}(\{\Psi\in\mathscr{L}(\mathscr{O})-\{0\}:\Psi<\Phi\}).
    \] 
\end{lemma}

\begin{observation}\label{obs:pseudospheres-regular-cw}
    If $\mathfrak{A}$ is a signed arrengement of pseudospheres in $S^{d-1}$ indexed by the finite set $V$ and with associated oriented matroid $\mathscr{O}$, then using Lemma \ref{lem:disks-in-oms} the arrangement $\mathfrak{A}$ gives a regular CW decomposition of $S^{d-1}$ whose open cells are of the form $B^-(\Phi^-)\cap S(\Phi^0)\cap B^+(\Phi^+)$. There is one cell for each $\Phi\in\mathscr{L}(\mathscr{O})-\{0\}$, and if $\mathfrak{A}$ is not essential then there are additional cells used to decompose $S(V)$.
\end{observation}

\begin{remark}\label{rmk:reorientation}
    Given a signed arrangement $\mathfrak{A}$ of pseudospheres in $S^{d-1}$ indexed by $V$ with associated oriented matroid $\mathscr{O}$, and a subset $U\subseteq V$, there is another signed arrangement $\mathfrak{A}_{-U}$ of pseudospheres in $S^{d-1}$ indexed by $V$, with associated oriented matroid denoted by $\mathscr{O}_{-U}$, for which $\mathfrak{s}_{\mathfrak{A}_{-U}}(\varphi)(v)$ is equal to $\mathfrak{s}_\mathfrak{A}(\varphi)(v)$ if $v\in V-U$ and $-\mathfrak{s}_\mathfrak{A}(\varphi)(v)$ if $v\in U$. Then for $v\in V$ and $\mathfrak{s}\in\{-,+\}$ the sets $B^\mathfrak{s}_v$ and $D^\mathfrak{s}_v$ defined in $\mathfrak{A}_{-U}$ are respectively equal to $B^\mathfrak{s}_v$ and $D^\mathfrak{s}_v$ defined in $\mathfrak{A}$ if $v\in V-U$, or to $B^{-\mathfrak{s}}_v$ and $D^{-\mathfrak{s}}_v$ defined in $\mathfrak{A}$ if $v\in U$. This can be combined with the last corollary to show that any intersection of sets of the form $B^{\mathfrak{s}}_v$ and $D^\mathfrak{s}_v$ is empty contractible, given at least one $B^\mathfrak{s}_v$ is used.
\end{remark}

\begin{observation}\label{obs:curly-vs-print}
    Let $\mathfrak{A}$ be a signed arrangement of pseudospheres in $S^{d-1}$ indexed by $V$, and $\mathscr{O}$ the associated oriented matroid. Let $\mathfrak{F}_\mathfrak{A}$ be the set of all formulas which can be constructed using the symbols $(-)\cap(-)$, $(-)\cup(-)$, as well as $S_v$, $B^\mathfrak{s}_v$, $D^\mathfrak{s}_v$ for $\mathfrak{s}\in\{-,+\}$ and $v\in V$. Similarly, define $\mathfrak{F}_\mathscr{O}$ as the set of formulas which can be constructed using $(-)\cap(-)$, $(-)\cup(-)$, and $\mathscr{S}_v$, $\mathscr{B}^\mathfrak{s}_v$, $\mathscr{D}^\mathfrak{s}_v$ with $\mathfrak{s}\in\{+,-\}$ and $v\in V$. Then there is a bijection between $\mathfrak{F}_\mathfrak{A}$ and $\mathfrak{F}_\mathscr{O}$. Let $\mathcal{F}$ be a formula in $\mathfrak{F}_\mathfrak{A}$, $F$ the subset of $S^{d-1}$ defined by $\mathcal{F}$, and let $\mathscr{F}$ be the subset of $\mathscr{L}(\mathscr{O})$ defined by the corresponding formula in $\mathfrak{F}_\mathscr{O}$. Then $\mathfrak{s}_\mathfrak{A}^{-1}(\mathscr{F})=F$, $\mathscr{F}\supseteq \mathfrak{s}_\mathfrak{A}(F)$, and $\mathscr{F}-\mathfrak{s}_\mathfrak{A}(F)$ is either $\emptyset$ or $\{0\}$. In particular, $F=\mathfrak{s}_\mathfrak{A}^{-1}(\mathscr{F})$ and $\mathfrak{s}_\mathfrak{A}(F)=\mathscr{F}$ hold under the condition that $\mathcal{F}$ is obtained using $(-)\cap(-)$, $(-)\cup(-)$, and formulas of the form $(B^\mathfrak{s}_v)\cap(\mathcal{F}')$, where $\mathfrak{s}\in\{-,+\}$, $v\in V$, and $\mathcal{F}'\in\mathfrak{F}_\mathfrak{A}$.

    To prove all of this, observe that by replacing $(-)\cap(-)$ with ``$(-)$ and $(-)$'', $(-)\cup(-)$ with ``$(-)$ or $(-)$'', $S_v$ with $\Xi(v)=0$, $B^\mathfrak{s}_v$ with $\Xi(v)=\mathfrak{s}$, and $D^\mathfrak{s}_v$ with $\Xi(v)\neq -\mathfrak{s}$, we get a formula $\mathcal{F}'(\Xi)$ with a free variable $\Xi$ such that $F=\{\varphi\in S^{d-1}:\mathcal{F}'(\mathfrak{s}_\mathfrak{A}(\varphi))\}$ and $\mathscr{F}=\{\Phi\in\mathscr{L}(\mathscr{O}):\mathcal{F}'(\Phi)\}$.
\end{observation}  

\begin{proposition}\label{prop:htpy-type-zero-avoiding}
    Let $\mathfrak{A}$ be a signed arrangement of pseudospheres in $S^{d-1}$ indexed by the finite set $V$, and let $\mathscr{O}$ be its associated oriented matroid. The link $\link_{\mathscr{C}_\mathscr{O}}(V-U)=(\mathscr{C}^*_{\mathscr{O}}[U])^*$ for $U\subseteq V$ can be described as follows:
    \begin{enumerate}[label=\normalfont{(\arabic*)}]
        \item\label{prop:htpy-type-zero-avoiding:void} If $\mathscr{B}^+(V-U)=\emptyset$ then $\link_{\mathscr{C}_{\mathscr{O}}}(V-U)$ is the void complex.
        \item\label{prop:htpy-type-zero-avoiding:nerve} If $\mathscr{B}^+(V-U)\neq\emptyset$ then $\link_{\mathscr{C}_{\mathscr{O}}}(V-U)$ is homotopy equivalent to $B^+\langle U\rangle\cap B^+(V-U)$.
    \end{enumerate}
    In particular, if we are in case \ref{prop:htpy-type-zero-avoiding:nerve} then $\link_{\mathscr{C}_\mathscr{O}}(V-U)$ is homotopy equivalent to an open subset of $S^{d-1}$. The assumption $\mathscr{B}^+(V-U)\neq\emptyset$ is automatically satisfied if $U=V$, in which case we have the implications:
    \begin{enumerate}[label=\normalfont{(2.\arabic*)}]
        \item\label{prop:htpy-type-zero-avoiding:sphere} If $\mathscr{D}^+(V)=\{0\}$ then $\mathscr{C}_\mathscr{O}=\link_{\mathscr{C}_\mathscr{O}}(\emptyset)\simeq S^{\operatorname{rank}(\mathscr{O})-1}$.
        \item\label{prop:htpy-type-zero-avoiding:entire-contr} Otherwise $\mathscr{C}_\mathscr{O}=\link_{\mathscr{C}_\mathscr{O}}(\emptyset)$ is contractible.
    \end{enumerate}
    If we are in case \ref{prop:htpy-type-zero-avoiding:nerve} but under the assumption that $U\subsetneq V$ then:
    \begin{enumerate}[resume,label=\normalfont{(2.\arabic*)}]
        \item\label{prop:htpy-type-zero-avoiding:non-Farkas} If $\mathscr{B}^+(V-U)\cap \mathscr{D}^-(U)=\emptyset$ then $\link_{\mathscr{C}_{\mathscr{O}}}(V-U)$ is contractible.
        \item\label{prop:htpy-type-zero-avoiding:boundary} Otherwise $\link_{\mathscr{C}_\mathscr{O}}(V-U)$ is also homotopy equivalent to
        \[
            S^{\operatorname{rank}(U')-1}*\left(S(U')\cap B^+\langle U-U'\rangle\cap B^+(V-U)\right)
        \]
        for any $U'\subseteq U$ with $\mathscr{D}^+(U')=\mathscr{S}(U')$, where by convention $S^{-1}=\emptyset$.
    \end{enumerate}
    %In particular, if we are in case \ref{prop:htpy-type-zero-avoiding:boundary} then $\link_{\mathscr{C}_\mathscr{O}}(V-U)$ is homotopy equivalent to an open subset of $S^{d-2}$.
\end{proposition}

\begin{proposition}\label{prop:htpy-type-elem-avoiding}
    Let $\mathfrak{A}$ be a signed arrangement of pseudospheres in $S^{d-1}$ indexed by the finite set $V\sqcup\{v_0\}$, and let $\mathscr{O}$ be its associated oriented matroid. The link $\link_{\mathscr{C}_{v_0,\mathscr{O}}}(V-U)=(\mathscr{C}_{v_0,\mathscr{O}}^*[U])^*$ can be desribed as follows:
    \begin{enumerate}[label=\normalfont{(\arabic*)}]
        \item\label{prop:htpy-type-elem-avoiding:void} If $\mathscr{D}^+(V-U)\cap\mathscr{B}^-_{v_0}=\emptyset$ then $\link_{\mathscr{C}_{v_0,\mathscr{O}}}(V-U)$ is the void complex.
        \item\label{prop:htpy-type-elem-avoiding:nerve} If $\mathscr{D}^+(V-U)\cap\mathscr{B}^-_{v_0}\neq\emptyset$ then $\link_{\mathscr{C}_{v_0,\mathscr{O}}}(V-U)\simeq D^+\langle U\rangle\cap D^+(V-U)\cap B^-_{v_0}$.
    \end{enumerate}
    In particular, if we are in case \ref{prop:htpy-type-elem-avoiding:nerve} then $\link_{\mathscr{C}_{v_0,\mathscr{O}}}(V-U)$ is homotopy equivalent to a subset of the open hemisphere $B^-_{v_0}$ of $S^{d-1}$. Moreover, if $v_0$ is not a loop in $\mathscr{O}$ then we have the implications:
    \begin{enumerate}[label=\normalfont{(2.\arabic*)}]
        \item\label{prop:htpy-type-elem-avoiding:nonpointed} If $\mathscr{B}^+(V\cup\{v_0\})=\emptyset$ then $\mathscr{C}_{v_0,\mathscr{O}}=\link_{\mathscr{C}_{v_0,\mathscr{O}}}(\emptyset)$ is contractible.
        \item\label{prop:htpy-type-elem-avoiding:pointed} Otherwise $\mathscr{C}_{v_0,\mathscr{O}}=\mathscr{C}_{\mathscr{O}/v_0}$.
    \end{enumerate}
    Therefore in point \ref{prop:htpy-type-elem-avoiding:pointed} the links in $\mathscr{C}_{v_0,\mathscr{O}}$ are described by Proposition \ref{prop:htpy-type-zero-avoiding}.
\end{proposition}

\begin{theorem}\label{thm:nerve-closed-cellular}
    Let $X$ be a regular CW complex, $A\subseteq X$ an intersection of a subcomplex and an open union of some open cells of $X$, and $(Z_i)_{i\in I}$ a finite collection of subcomplexes of $X$ with $A\subseteq\bigcup_{i\in I}Z_i$. If for all $\emptyset\neq J\subseteq I$ the space $A\cap\bigcap_{j\in J}Z_j$ is empty or contractible, then $A$ is homotopy equivalent to the nerve $\mathscr{N}_{\mathfrak{Z}}$ of the cover $\mathfrak{Z}:=(A\cap Z_i)_{i\in I}$ of $A$.
\end{theorem}

\begin{proof}[Proof of Proposition \ref{prop:htpy-type-elem-avoiding}]
    Consider $\mathscr{C}_{v_0,\mathscr{O}}$ as the simplicial complex given by $\{U\subseteq V:\mathscr{D}^+(U)\cap\mathscr{B}^-_{v_0}\neq\emptyset\}$. In particular, under the assumption of point \ref{prop:htpy-type-elem-avoiding:void} the set $V-U$ is not a face of $\mathscr{C}_{v_0,\mathscr{O}}$, so its link is the void complex:
    \[
        \link_{\mathscr{C}_{v_0,\mathscr{O}}}(V-U)=\{U'\in\mathscr{C}_{v_0,\mathscr{O}}:(V-U)\cap U'=\emptyset,(V-U)\cup U'\in\mathscr{C}_{v_0,\mathscr{O}}\},
    \]
    but $(V-U)\cup U'\in\mathscr{C}_{v_0,\mathscr{O}}\implies V-U\in\mathscr{C}_{v_0,\mathscr{O}}$, so the latter set is empty, proving point \ref{prop:htpy-type-elem-avoiding:void}.

    For point \ref{prop:htpy-type-elem-avoiding:nerve}, consider the collection
    \[
        \mathfrak{U}:=(D^+_u\cap D^+(V-U)\cap B^-_{v_0})_{u\in U}
    \]
    of subsets of $D^+(V-U)\cap B^-_{v_0}$, which form a cover of $D^+\langle U\rangle\cap D^+(V-U)\cap B^-_{v_0}$. If $\emptyset\neq U'\subseteq U$, then $\bigcap_{u'\in U'}(D^+_{u'}\cap D^+(V-U)\cap B^-_{v_0})=D^+(V-(U-U'))\cap B^-_{v_0}$ is empty or contractible by the end of Remark \ref{rmk:reorientation}. Note that this intersection is nonempty if and only if $\mathscr{D}^+(V-(U-U'))\cap\mathscr{B}^-_{v_0}\neq\emptyset$ (Observation \ref{obs:curly-vs-print}), or equivalently if $V-(U-U')\in\mathscr{C}_{v_0,\mathscr{O}}$ by the description of $\mathscr{C}_{v_0,\mathscr{O}}$, but this holds if and only if $U-U'\notin\mathscr{C}_{v_0,\mathscr{O}}^*$, which is the same as saying $U'\in(\mathscr{C}_{v_0,\mathscr{O}}^*[U])^*$. In particular, to see that the nerve $\mathscr{N}_{\mathfrak{U}}$ of the cover $\mathfrak{U}$ is just $(\mathscr{C}_{v_0,\mathscr{O}}^*[U])^*=\link_{\mathscr{C}_{v_0,\mathscr{O}}}(V-U)$ it suffices to check that $\emptyset\in\mathscr{N}_\mathfrak{U}\iff\emptyset\in\link_{\mathscr{C}_{v_0,\mathscr{O}}}(V-U)$, but nerves are by definition not void and $\mathscr{D}^+(V-U)\cap\mathscr{B}^-_{v_0}\neq\emptyset\implies V-U\in\mathscr{C}_{v_0,\mathscr{O}}\implies\emptyset\in\link_{\mathscr{C}_{v_0,\mathscr{O}}}(V-U)$. Therefore, the only thing needed to prove point \ref{prop:htpy-type-elem-avoiding:nerve} is to verify that the assumptions of the Nerve Theorem (Theorem \ref{thm:nerve-closed-cellular}) hold. 
    Let $X:=S^{d-1}$ with the regular cell decomposition given by $\mathfrak{A}$ (see Observation \ref{obs:pseudospheres-regular-cw}), $A:=D^+\langle U\rangle\cap D^+(V-U)\cap B^-_{v_0}$, and $Z_u:=D^+_u$ for $u\in U$. Then $A\subseteq\bigcup_{u\in U}Z_u$, $Z_u$ is easily seen to be a subcomplex of $X$, $A$ is indeed a union of open cells, and the nerve of the cover $(A\cap Z_u)_{u\in U}$ of $A$ is equal to $\mathscr{N}_\mathfrak{U}$, so Theorem \ref{thm:nerve-closed-cellular} is applicable and gives that $\mathscr{N}_\mathfrak{U}=\mathscr{N}_{(A\cap Z_u)_{u\in U}}\simeq A=D^+\langle U\rangle\cap D^+(V-U)\cap B^-_{v_0}$.

    For point \ref{prop:htpy-type-elem-avoiding:nonpointed} note that when $U:=V$ the condition $\mathscr{D}^+(V-U)\cap\mathscr{B}^-_{v_0}\neq\emptyset$ is automatically satisfied as $v_0$ is not a loop, so 
    \begin{multline*}
    \mathscr{C}_{v_0,\mathscr{O}}=\link_{\mathscr{C}_{v_0,\mathscr{O}}}(\emptyset)\simeq D^+\langle V\rangle\cap D^+(\emptyset)\cap B^-_{v_0}=D^+\langle V\rangle\cap B^-_{v_0}\\ =(S^{d-1}-B^-(V))\cap B^-_{v_0}=B^-_{v_0}-B^-(V\cup\{v_0\}).	
    \end{multline*}
    However, $B^-(V\cup\{v_0\})$ is empty because $\mathscr{B}^-(V\cup\{v_0\})$ is empty (Observation \ref{obs:curly-vs-print}), which holds as $\mathscr{B}^+(V\cup\{v_0\})$ is empty, so in this case $\mathscr{C}_{v_0,\mathscr{O}}\simeq B^-_{v_0}-B^-(V\cup\{v_0\})=B^-_{v_0}$, but this space is homeomorphic to an open hemisphere of $S^{d-1}$ because $v_0$ is not a loop, and hence is contractible.

    It is left to prove point \ref{prop:htpy-type-elem-avoiding:pointed}, namely that $\mathscr{C}_{v_0,\mathscr{O}}=\mathscr{C}_{\mathscr{O}/v_0}$ if $\mathscr{B}^+(V\cup\{v_0\})\neq\emptyset$.
    After expanding the definitions of both complexes, we just have to show that for any $\emptyset\neq U\subseteq V$ we have that $\mathscr{D}^+(U)\cap\mathscr{B}^-_{v_0}\neq\emptyset$ in $\mathscr{O}$ if and only if $\mathscr{B}^+(U)\neq\emptyset$ in $\mathscr{O}/v_0$, but this latter condition is equivalent to saying that $\mathscr{B}^+(U)\cap\mathscr{S}_{v_0}\neq\emptyset$ in $\mathscr{O}$. Hence from this point on we work in $\mathscr{O}$ and aim to show $\mathscr{D}^+(U)\cap\mathscr{B}^-_{v_0}\neq\emptyset\iff\mathscr{B}^+(U)\cap\mathscr{S}_{v_0}\neq\emptyset$.
    By assumption we have a $\Phi\in\mathscr{B}^+(V\cup\{v_0\})$. If $\Psi\in\mathscr{B}^+(U)\cap\mathscr{S}_{v_0}$ then $\Psi\circ(-\Phi)\in\mathscr{B}^+(U)\cap\mathscr{B}^-_{v_0}\subseteq\mathscr{D}^+(U)\cap\mathscr{B}^-_{v_0}$. For the other direction, first take a $\Psi_0\in\mathscr{D}^+(U)\cap\mathscr{B}^-_{v_0}$, then pick a $\Theta$ given by axiom \ref{def:om:elim} for $\Pi:=\Phi$, $\Xi:=\Psi_0\circ\Phi$, and $v:=v_0$. For $u\in U$ both $\Pi(u)=+$ and $\Xi(u)=\Psi_0(u)\circ+=+$, so $\Theta(u)=+$ too, but also $\Theta(v_0)=0$, so $\Theta\in\mathscr{B}^+(U)\cap\mathscr{S}_{v_0}$.
\end{proof}

\end{document}
